\section{Metodología de Implementación}

\subsection{Enfoque metodológico del proyecto}

La implementación del sistema Odoo 19 en \textbf{Machu Picchu Exclusive Tours E.I.R.L.} se diseñó como un proceso de \textbf{transformación operativa incremental}, priorizando el uso de funcionalidades estándar de la plataforma para minimizar riesgos técnicos y asegurar una curva de adopción razonable para el equipo.

El enfoque se centró en:

\begin{itemize}
    \item Entregar primero la base del CRM y la estructura de datos maestros.
    \item Digitalizar progresivamente el flujo operativo de reservas.
    \item Consolidar en una etapa final el control de saldos, la capacitación y las pruebas integrales.
\end{itemize}

\subsection{Metodología ágil aplicada}

El proyecto siguió una adaptación de la metodología \textbf{ágil} basada en sprints, con revisiones periódicas con la dirección de la agencia:

\begin{itemize}
    \item \textbf{Entregas parciales por módulos:} Primero CRM y configuración de instancia; luego tablero de proyectos; finalmente, facturación básica y capacitación.
    \item \textbf{Retroalimentación continua (Feedback):} Ajustes en las etapas del tablero y en los campos de CRM conforme se revisaban casos reales de clientes.
    \item \textbf{Flexibilidad en el alcance:} Manteniendo las exclusiones definidas, pero adaptando nomenclaturas y vistas para ajustarlas al lenguaje de la operación turística.
\end{itemize}

\subsection{Fases de implementación}

Las fases prácticas del proyecto se alinearon con los sprints:

\begin{table}[h]
\centering
\renewcommand{\arraystretch}{1.5}
\begin{tabular}{|l|p{4cm}|p{6cm}|p{3cm}|}
\hline
\textbf{Fase} & \textbf{Módulo / Actividad Clave} & \textbf{Actividad Principal} & \textbf{Situación} \\ \hline
Fase 1 & Diagnóstico y Planificación & Validación de flujos comerciales y operativos, mapeo de roles y definición de requerimientos. & Finalizado \\ \hline
Fase 2 & Core \& CRM & Configuración de instancia en la nube, usuarios, permisos y fichas de PAX con adjuntos. & Finalizado \\ \hline
Fase 3 & Operaciones & Configuración del tablero de Proyectos con etapas Borrador, Confirmación, Compra y Ejecución. & Finalizado \\ \hline
Fase 4 & Actividades \& Alertas & Parametrización de actividades, recordatorios y vista de calendario. & Finalizado \\ \hline
Fase 5 & Finanzas \& QA & Configuración de facturación básica, control de adelantos y saldos, capacitaciones y pruebas integrales. & Finalizado \\ \hline
\end{tabular}
\caption{Fases de implementación del proyecto en Machu Picchu Exclusive Tours.}
\end{table}

